{ % Tabellendefinitionen gelten nur in diesem Dokument

% Longtable auf Spaltebreite ausdehnen
\setlength\LTleft{0pt}
\setlength\LTright{0pt}
% Definition neuer Spaltentypen
\newcolumntype{Z}{>{$}p{0.15\linewidth}<{$}} % Symbol (mit Matheumgebung)
\newcolumntype{A}{>{\bfseries} p{0.15\linewidth}} % Abkürzung (fett)
\newcolumntype{B}{l @{\extracolsep{\fill}}} % Beschreibung


\chapter{Symbolverzeichnis}

\section*{Formelzeichen}

\begin{longtable}{Z B r}
\textbf{Symbol} & \textbf{Bedeutung}         & \textbf{Einheit} \\
\addlinespace
\endhead
% Symbolspalte wird immer in Matheumgebung gesetzt --> $ $ nicht nötig
%************************************************************************
a               & Länge                      & \SI{}{m}         \\
b               & Breite                     & \SI{}{m}         \\
c               & spezifische Wärmekapazität & \SI{}{kJ/(kg.K)} \\
\ell            & charakteritische Länge     & \SI{}{m}         \\
Re              & Reynoldszahl               &                  \\ 
w               & Strömungsgeschwindigkeit   & \SI{}{m/s}       \\ 
%************************************************************************
\end{longtable}


\section*{Griechische Symbole}

\begin{longtable}{Z B r}
\textbf{Symbol} & \textbf{Bedeutung}         & \textbf{Einheit} \\
\addlinespace
\endhead
% Symbolspalte wird immer in Matheumgebung gesetzt --> $ $ nicht nötig
%************************************************************************
\lambda         & Wärmeleitfähigkeit         & \SI{}{W/(m~K)}   \\ 
\nu             & kinematische viskosität    & \SI{}{m^2/s}     \\
\varrho         & Dichte                     & \SI{}{kg/m^3}     \\
%************************************************************************
\end{longtable}


\section*{Indizes}

\begin{longtable}{Z l} 
\textbf{Symbol} & \textbf{Bedeutung} \\
\addlinespace
\endhead
% Symbolspalte wird immer in Matheumgebung gesetzt --> $ $ nicht nötig
%************************************************************************
\mathrm{L}      & Luft               \\
\mathrm{p}      & isobar             \\
%************************************************************************
\end{longtable}



\chapter{Abkürzungsverzeichnis}

\begin{longtable}{@{} A l @{}}
%************************************************************************
eps & Encapsulated PostScript  \\
pdf & Portable Document Format \\
svg & Scalable Vector Format   \\
%************************************************************************
\end{longtable}
}


\setcounter{table}{0}
% setzt den Tabellenzähler zurück, damit das Symbolverzeichnis
% nicht mitgezählt wird


%%% Local Variables: 
%%% mode: latex
%%% TeX-master: "master"
%%% coding: utf-8
%%% End: 